% Sec. IV Time and Space Alignment (about 0.8 page).
\section{Time and Space Alignment}
\label{sec:alignment}
This section answers how well the testbed puts every measurement from
every agent on one clock and in one rig frame. Each quantity is
validated with data that was not used in its estimation.

\subsection{Clock}
\label{sec:clock}
The clocks are synchronized with chrony. \texttt{mobile\_1} acts as
the server, and every other agent and \texttt{sniffer} act as clients
polling every \SI{8}{\second}. This traffic passes over the measured
link, but at a few hundred bytes per poll it is negligible, and since
it is present in every run it forms part of the measured channel.
Recording starts only after every unit reports a converged offset.
Each unit logs its offset and any clock adjustment
(\texttt{ntp/status}, \texttt{ntp/events}) into its own bag, so the
alignment at any frame can be checked and poorly aligned frames can
be filtered out. Two numbers describe the result, and they answer
different questions. The RMS offset chrony reports per client is tens
of microseconds (median \SIrange{10.9}{42.5}{\micro\second} across
clients on a live pass, no clock stepped; supplementary Table S-II).
The \emph{measured} offset over the whole of the coop\_2 run, the
quantity a user of the data cares about, is at most
\SI{0.99}{\milli\second}, and chrony's strict error bound, the
offset plus half the root delay, is at most \SI{12}{\milli\second}.
We report the measured offset as the clock guarantee and carry the
strict bound through the error budget as the worst case; the clock
rows of Table~\ref{tab:calib} give both. \todo{Repeat over every
pass; report the worst pairwise value.}

\subsection{Sniffer Bias}
\label{sec:sniffer}
These offsets describe the stability of each clock relative to the
server, not its accuracy. NTP assumes a symmetric path delay, and any
asymmetry appears as a constant bias that chrony cannot observe. This
is relevant only for \texttt{sniffer}, which reaches the server over
its \SI{2.4}{\giga\hertz} management link with a path delay several
times larger than that of the other units, so its bias may be as large
as half that delay, \SI{7.8}{\milli\second}. We measure the bias
from the probe packets that \texttt{mobile\_2} sends and
\texttt{sniffer} captures. Each packet is stamped by the sending
application on the clock of \texttt{mobile\_2} and on capture on the
clock of \texttt{sniffer}; the difference is the bias plus the time the
packet spent queued and in the drivers, which grows with load. The bias
is therefore estimated from the lower envelope of the differences, their
minimum over windows in which the link is idle \cite{paxson1998calibrating},
which bounds bias plus the irreducible driver latency: \todo{value}.

\subsection{Error Budget}
\label{sec:budget}
Agreeing clocks are not sufficient if the stamps themselves are late.
Every message is stamped on arrival at its host, so each stamp lags
its measurement by the transport latency of that sensor
\todo{estimates per sensor: LiDAR over Ethernet, cameras over USB,
radar over UART}, and camera stamps lag further by the exposure time.
These lags are approximately constant for a given sensor; we leave the
stamps as recorded so that users can apply their own correction.

Every timing error $\Delta t$ displaces a measurement taken on a
platform moving at speed $v$ by
\begin{equation}
  \delta x = v \, \Delta t .
  \label{eq:time_bound}
\end{equation}
We take $v$ as the measured 95th-percentile agent speed,
\todo{$v_{95}$} on \texttt{coop\_2}; at \SI{1}{\metre\per\second} the
measured clock offset of \SI{0.99}{\milli\second} corresponds to
\SI{1}{\milli\metre}, and the strict \SI{12}{\milli\second} bound to
\SI{1.2}{\centi\metre}. The sniffer bias of \todo{} corresponds
to \todo{}, which is \todo{}\,\% of the half-wavelength
$\lambda/2 = \SI{26.1}{\milli\metre}$ at \SI{5745}{\mega\hertz} that
sets the CSI sampling requirement (Sec.~\ref{sec:csi}). Transport
latency and exposure are larger: a constant lag of
\SI{10}{\milli\second} corresponds to \SI{3}{\centi\metre}, but since
it is constant it can be estimated and removed. The largest term is
the LiDAR sweep, \SI{100}{\milli\second} at \SI{10}{\hertz}, so the
first and last points of a scan can be \SI{30}{\centi\metre} apart;
it is removed by per-point deskewing with IMU-propagated motion in the
reference pipeline \cite{glim}. Combining the terms that remain after
compensation, the worst-case intra-frame spatial error at $v_{95}$ is
\begin{equation}
  \delta x_{\max} = v_{95}\,\big(\Delta t_{\mathrm{clk}}
  + \Delta t_{\mathrm{lat}} + \Delta t_{\mathrm{exp}}\big)
  = \todo{Y}\,\si{\centi\metre},
  \label{eq:fusion_bound}
\end{equation}
where $\Delta t_{\mathrm{clk}}$ is the worst clock offset,
$\Delta t_{\mathrm{lat}}$ the largest uncorrected transport latency,
and $\Delta t_{\mathrm{exp}}$ the longest camera exposure. The budget
is dominated by sensor physics and transport, not by synchronization:
measured clock alignment contributes a few millimetres at most, while
the larger terms are either compensated or constant and removable. \todo{Show the
budget as a stacked bar; it feeds the headroom figure of E1.}

\subsection{Extrinsics}
\label{sec:spatial}
Three transforms are estimated per platform. On \texttt{mobile\_1},
\Ttransform{Base}{LiDAR} is obtained by hand-eye calibration: the
robot is driven through a sequence of motions, and the transform is
solved from the pairs of relative motions reported by wheel odometry
and by LiDAR odometry \cite{furrer2018handeye}. On both mobile agents,
\Ttransform{LiDAR}{Camera} is obtained by direct, target-free alignment
of image and point-cloud features \cite{koide2023lidarcam}. For the
radar, whose returns are too sparse for feature alignment,
\Ttransform{LiDAR}{Radar} is obtained with a corner reflector placed at
several positions \cite{domhof2019calib}. For the infrastructure rig,
the camera and radar are mounted on the same fixed bracket and each is
calibrated against a LiDAR by the same methods, which gives
\Ttransform{Camera}{Radar} directly; the world pose of each mast camera
is obtained through the fiducial chain of Sec.~\ref{sec:fiducials}.

Each extrinsic is checked on held-out data. For
\Ttransform{LiDAR}{Camera}, LiDAR returns on the ChArUco boards of the
mapping session are projected into the image and compared with the
detected corners; for the RGB-D agent, whose trajectory uses board
sightings as constraints, the boards are split once per site into an
estimation set and a held-out set, and the residual is evaluated on
the held-out set only \todo{state the split}. The board detector
itself is validated before it validates anything: re-run on the
mapping session's images it reproduces the pipeline's board poses to
\SI{8.1}{\milli\metre} / 0.66$^\circ$ on the ZED and
\SI{8.3}{\milli\metre} / 0.58$^\circ$ on the RealSense, at or below
each board's own scatter (\SI{2.0}{} and \SI{6.9}{\milli\metre}),
with lens distortion handled at corner interpolation rather than at
PnP, which otherwise costs \SI{34}{\milli\metre}. For
\Ttransform{LiDAR}{Radar} and \Ttransform{Camera}{Radar}, the reflector
is placed at positions not used in calibration and the radar detection
is compared with the LiDAR centroid. For \Ttransform{Base}{LiDAR}, the
hand-eye result is checked through the board residual of
Sec.~\ref{sec:gt_val}, which tests \Ttransform{Base}{LiDAR} and
\Ttransform{LiDAR}{Camera} at once. Table~\ref{tab:calib} lists the
residuals; the same values are stored with the data
(\texttt{calibration/residuals.json}) so that users can weight each
transform by its accuracy.

\begin{table}[t]
\centering
\caption{Calibration residuals on held-out observations.}
\label{tab:calib}
\footnotesize
\setlength{\tabcolsep}{4pt}
\begin{tabular}{@{}llll@{}}
\toprule
Transform & Check & Median & p95 \\
\midrule
\Ttransform{LiDAR}{Camera}  & board corner reprojection (px) & \todo{} & \todo{} \\
Board detector, ZED / RealSense & board pose vs.\ mapping pipeline (mm, deg) & 8.1, 0.66 / 8.3, 0.58 & \todo{} \\
\Ttransform{LiDAR}{Radar}   & reflector offset (cm)          & \todo{} & \todo{} \\
\Ttransform{Base}{LiDAR}    & board pose via chain (cm, deg) & \todo{} & \todo{} \\
\Ttransform{Camera}{Radar}  & reflector offset (cm), infra   & \todo{} & \todo{} \\
Clock, all clients          & chrony RMS offset to \texttt{mobile\_1} (\si{\micro\second}) & 10.9--42.5 & \todo{} \\
Clock, coop\_2 run           & measured offset / strict bound (ms) & $\leq 0.99$ / $\leq 12$ & -- \\
Clock, \texttt{sniffer}     & probe-packet bias (ms)         & \todo{} & \todo{} \\
\bottomrule
\end{tabular}
\end{table}

